\section{Discussion}
\subsection{Interpreting apparently conflicting results}
The evidence does not constitute a four-way controlled contest. REEF and Tally protect priority inference using different interception/preemption paths and different comparators; Tally's background-throughput disadvantage at one idle-time point does not refute REEF's p99 result in a GV100 PyTorch mix \cite{reef,tally}. It instead shows that a latency objective and best-effort throughput can conflict. REEF also reports a separate TensorRT comparison that weakens any unconditional throughput claim for its own design \cite{reef}. The responsible synthesis names the engine, arrival process, percentile and denominator for each finding, and evaluates whether the objective was a hard priority SLO or total throughput.

MPS packing and MPS interference are compatible observations. VastPipe's measured throughput in an A100 inference pair suggests room for quota optimization, while DeInfer's RTX 3080 experiment shows residual latency interference under MPS limits \cite{vastpipe,deinfer}. Their devices, models, scheduler logic and reported statistics differ. The SoCC Green Context characterization provides a mechanistic explanation for some remaining contention: separated SM sets do not separate shared cache and off-chip bandwidth in its H100 experiments \cite{interference_deeper}. This is a plausible explanation across studies, explicitly an inference rather than a direct cross-paper causal test.

Hybrid MIG/MPS outcomes similarly depend on the partition profile and model fit. BOER reports a local adverse large-model case against MIG-only serving, while other hybrid systems optimize different objective functions and often report server-wide gains \cite{boer,parva,miger}. GNN-Triton's thrashing and Dyna-P's JCT trade-off show how profile size and placement can change the answer even before the next scheduling layer is added \cite{gnntriton,dynap}. MIGServe's four-GPU SLO-attainment rate and Protean's serverless evaluation should inform system design, but cannot be used to assign a single-device causal multiplier to MIG \cite{migserve,protean}.

\subsection{Evidence strength and transfer}
The strongest claims in this review are bounded within a study: a specified source location, hardware, workload and comparator support the stated result. Transferring the direction to an unseen GPU generation, model or arrival trace is weaker. The 43 current quality scores total 355 points (mean 8.26); 36 are high and seven moderate under the protocol's reporting rubric. These scores do not imply that 36 studies directly compare all four mechanisms. Only five studies receive a full point for measurement controls (Q7), while 35 receive 0.5 and three receive zero. Nine receive a full point for comparator adequacy (Q5), while 34 receive 0.5. These item patterns are more informative for engineering confidence than the mean score alone; they follow the item matrix in Appendix B.

Dropping the seven moderate studies leaves 36 current high-score studies. That changes which observations support serverless sharing, fairness or dynamic placement: VastPipe, ArkGPU, Dyna-P, InferFair, Deft, FaST-GShare and Edge-sky are among the moderate rows in the current quality sheet. It does not make the surviving results directly comparable, and it cannot answer whether the still-unassessed reports contain contrary evidence. \red{Repeat this sensitivity analysis against the final included set and the final extraction, not just the 43-study checkpoint.}

\subsection{From evidence to an operational decision}
The decision starts with non-negotiable constraints. A strict p99 service with heterogeneous best-effort co-tenants is not adequately specified by a target average utilization. A deployment with a hard memory partition or fault-containment requirement is not adequately specified by lower measured interference. A server with an unsupported MIG GPU cannot choose a MIG profile, however favorable a different hardware study appears. Once those constraints are written, choose a representative model and request trace, measure solo latency, then vary tenant count and mechanism under the same SLO and resource budget. Record tail distributions, SLO goodput, transition/reconfiguration costs and the particular shared resources stressed. The current evidence supplies candidates and known failure modes for that test; it does not replace it.

The review's contribution is a decision boundary rather than a global leaderboard. Temporal systems have the most direct route to reserving or reclaiming execution time for a priority request but may sacrifice background work. MPS can fill unused compute with modest partition management but does not remove all contention. MIG offers explicit instances where hardware and memory profile fit, yet partition granularity and switching costs can waste capacity. API remoting can centralize control at the cost of another data and trust path. Hybrids inherit more knobs and a larger configuration-search burden. The citations in Sections 5--8 support these qualitative contrasts; \red{their final ranking under any particular deployment remains conditional on completing the review and measuring that deployment.}
